\section{AHSME 1975}
        \begin{enumerate}
        	\item (\textit{AHSME 1975}) The value of $\frac {1}{2 - \frac {1}{2 - \frac {1}{2 - \frac12}}}$ is 


 $\textbf{(A)}\ 3/4 \qquad  \textbf{(B)}\ 4/5 \qquad  \textbf{(C)}\ 5/6 \qquad  \textbf{(D)}\ 6/7 \qquad  \textbf{(E)}\ 6/5$


 

	\item (\textit{AHSME 1975}) For which real values of m are the simultaneous equations 

 \begin{align*}y &= mx + 3 \\  y& = (2m - 1)x + 4\end{align*}
 
satisfied by at least one pair of real numbers $(x,y)$? 


 $\textbf{(A)}\ \text{all }m\qquad \textbf{(B)}\ \text{all }m\neq 0\qquad \textbf{(C)}\ \text{all }m\neq 1/2\qquad \textbf{(D)}\ \text{all }m\neq 1\qquad \\ \textbf{(E)}\ \text{no values of }m$


 

	\item (\textit{AHSME 1975}) Which of the following inequalities are satisfied for all real numbers $a, b, c, x, y, z$ which satisfy the conditions $x < a, y < b$, and $z < c$? 


 $\text{I}. \ xy + yz + zx < ab + bc + ca \\ \text{II}. \ x^2 + y^2 + z^2 < a^2 + b^2 + c^2 \\ \text{III}. \ xyz < abc$


 $\textbf{(A)}\ \text{None are satisfied.} \qquad  \textbf{(B)}\ \text{I only} \qquad  \textbf{(C)}\ \text{II only} \qquad  \textbf{(D)}\ \text{III only} \qquad  \textbf{(E)}\ \text{All are satisfied.}$


 

	\item (\textit{AHSME 1975}) If the side of one square is the diagonal of a second square, what is the ratio of the area of the first square to the area of the second? 


 $\textbf{(A)}\ 2 \qquad  \textbf{(B)}\ \sqrt2 \qquad  \textbf{(C)}\ 1/2 \qquad  \textbf{(D)}\ 2\sqrt2 \qquad  \textbf{(E)}\ 4$


 

	\item (\textit{AHSME 1975}) The polynomial $(x+y)^9$ is expanded in decreasing powers of $x$. The second and third terms have equal values 
when evaluated at $x=p$ and $y=q$, where $p$ and $q$ are positive numbers whose sum is one. What is the value of $p$? 


 $\textbf{(A)}\ 1/5 \qquad  \textbf{(B)}\ 4/5 \qquad  \textbf{(C)}\ 1/4 \qquad  \textbf{(D)}\ 3/4 \qquad  \textbf{(E)}\ 8/9$


 

	\item (\textit{AHSME 1975}) The sum of the first eighty positive odd integers subtracted from the sum of the first eighty positive even integers is 


 $\textbf{(A)}\ 0 \qquad  \textbf{(B)}\ 20 \qquad  \textbf{(C)}\ 40 \qquad  \textbf{(D)}\ 60 \qquad  \textbf{(E)}\ 80$


 

	\item (\textit{AHSME 1975}) For which non-zero real numbers $x$ is $\frac{|x-|x|\-|}{x}$ a positive integer? 


 $\textbf{(A)}\ \text{for negative } x \text{ only} \qquad \\ \textbf{(B)}\ \text{for positive } x \text{ only} \qquad \\ \textbf{(C)}\ \text{only for } x \text{ an even integer} \qquad \\ \textbf{(D)}\ \text{for all non-zero real numbers } x \\ \textbf{(E)}\ \text{for no non-zero real numbers } x$


 

	\item (\textit{AHSME 1975}) If the statement "All shirts in this store are on sale." is false, then which of the following statements must be true? 


 I. All shirts in this store are at non-sale prices. 


 II. There is some shirt in this store not on sale. 


 III. No shirt in this store is on sale. 


 IV. Not all shirts in this store are on sale. 


 

$\textbf{(A)}\ \text{II only} \qquad  \textbf{(B)}\ \text{IV only} \qquad  \textbf{(C)}\ \text{I and III only} \qquad  \textbf{(D)}\ \text{II and IV only} \qquad  \textbf{(E)}\ \text{I, II and IV only}$


 

	\item (\textit{AHSME 1975}) Let $a_1, a_2, \ldots$ and $b_1, b_2, \ldots$ be arithmetic progressions such that $a_1 = 25, b_1 = 75$, and $a_{100} + b_{100} = 100$. 
Find the sum of the first hundred terms of the progression $a_1 + b_1, a_2 + b_2, \ldots$


 $\textbf{(A)}\ 0 \qquad  \textbf{(B)}\ 100 \qquad  \textbf{(C)}\ 10,000 \qquad  \textbf{(D)}\ 505,000 \qquad \\ \textbf{(E)}\ \text{not enough information given to solve the problem}$


 

	\item (\textit{AHSME 1975}) The sum of the digits in base ten of $(10^{4n^2+8}+1)^2$, where $n$ is a positive integer, is 


 $\textbf{(A)}\ 4 \qquad  \textbf{(B)}\ 4n \qquad  \textbf{(C)}\ 2+2n \qquad  \textbf{(D)}\ 4n^2 \qquad  \textbf{(E)}\ n^2+n+2$


 

	\item (\textit{AHSME 1975}) Let $P$ be an interior point of circle $K$ other than the center of $K$. Form all chords of $K$ which pass through $P$, 
and determine their midpoints. The locus of these midpoints is 


 $\textbf{(A)}\ \text{a circle with one point deleted} \qquad \\ \textbf{(B)}\ \text{a circle if the distance from } P \text{ to the center of } K \text{ is less than} \\ \text{one half the radius of } K \text{; otherwise a circular arc of less than} 360^{\circ}\qquad \\ \textbf{(C)}\ \text{a semicircle with one point deleted} \qquad \\ \textbf{(D)}\ \text{a semicircle} \qquad  \textbf{(E)}\ \text{a circle}$


 

	\item (\textit{AHSME 1975}) If $a \neq b, a^3 - b^3 = 19x^3$, and $a-b = x$, which of the following conclusions is correct? 


 $\textbf{(A)}\ a=3x \qquad  \textbf{(B)}\ a=3x \text{ or } a = -2x \qquad  \textbf{(C)}\ a=-3x \text{ or } a = 2x \qquad \\ \textbf{(D)}\ a=3x \text{ or } a=2x \qquad  \textbf{(E)}\ a=2x$


 

	\item (\textit{AHSME 1975}) The equation $x^6 - 3x^5 - 6x^3 - x + 8$ has 


 $\textbf{(A)}\ \text{no real roots} \qquad  \textbf{(B)}\ \text{exactly two distinct negative roots} \qquad \\ \textbf{(C)}\ \text{exactly one negative root} \qquad  \textbf{(D)}\ \text{no negative roots, but at least one positive root} \qquad \\ \textbf{(E)}\ \text{none of these}$


 

	\item (\textit{AHSME 1975}) If the $whatsis$ is $so$ when the $whosis$ is $is$ and the $so$ and $so$ is $is \cdot so$, 
what is the $whosis \cdot whatsis$ when the $whosis$ is $so$, the $so$ and $so$ is $so \cdot so$, and the $is$ is two 
($whatsis, whosis, is$, and $so$ are variables taking positive values)? 


 $\textbf{(A)}\ whosis \cdot is \cdot so \qquad  \textbf{(B)}\ whosis \qquad  \textbf{(C)}\ is \qquad  \textbf{(D)}\ so\qquad \textbf{(E)}\ so\text{ and }so$


 

	\item (\textit{AHSME 1975}) In the sequence of numbers $1, 3, 2, \ldots$ each term after the first two is equal to the term preceding it minus the term preceding that. 
The sum of the first one hundred terms of the sequence is 


 $\textbf{(A)}\ 5 \qquad  \textbf{(B)}\ 4 \qquad  \textbf{(C)}\ 2 \qquad  \textbf{(D)}\ 1 \qquad  \textbf{(E)}\ -1$


 

	\item (\textit{AHSME 1975}) If the first term of an infinite geometric series is a positive integer, the common ratio is the reciprocal of a positive integer, 
and the sum of the series is $3$, then the sum of the first two terms of the series is 


 $\textbf{(A)}\ 1/3 \qquad  \textbf{(B)}\ 2/3 \qquad  \textbf{(C)}\ 8/3 \qquad  \textbf{(D)}\ 2 \qquad  \textbf{(E)}\ 9/2$


 

	\item (\textit{AHSME 1975}) A man can commute either by train or by bus. If he goes to work on the train in the morning, 
he comes home on the bus in the afternoon; and if he comes home in the afternoon on the train, 
he took the bus in the morning. During a total of $x$ working days, the man took the bus to work 
in the morning $8$ times, came home by bus in the afternoon $15$ times, and commuted by train 
(either morning or afternoon) $9$ times. Find $x$. 


 $\textbf{(A)}\ 19 \qquad  \textbf{(B)}\ 18 \qquad  \textbf{(C)}\ 17 \qquad  \textbf{(D)}\ 16 \qquad \\ \textbf{(E)}\ \text{not enough information given to solve the problem}$


 

	\item (\textit{AHSME 1975}) A positive integer $N$ with three digits in its base ten representation is chosen at random, 
with each three digit number having an equal chance of being chosen. The probability that $\log_2 N$ is an integer is 


 $\textbf{(A)}\ 0 \qquad  \textbf{(B)}\ 3/899 \qquad  \textbf{(C)}\ 1/225 \qquad  \textbf{(D)}\ 1/300 \qquad  \textbf{(E)}\ 1/450$


 

	\item (\textit{AHSME 1975}) Which positive numbers $x$ satisfy the equation $(\log_3x)(\log_x5)=\log_35$? 


 $\textbf{(A)}\ 3 \text{ and } 5 \text{ only} \qquad  \textbf{(B)}\ 3, 5, \text{ and } 15 \text{ only} \qquad \\ \textbf{(C)}\ \text{only numbers of the form } 5^n \cdot 3^m, \text{ where } n \text{ and } m \text{ are positive integers} \qquad \\ \textbf{(D)}\ \text{all positive } x \neq 1 \qquad  \textbf{(E)}\ \text{none of these}$


 

	\item (\textit{AHSME 1975}) In the adjoining figure $\triangle ABC$ is such that $AB = 4$ and $AC = 8$. If $M$ is the midpoint of $BC$ and $AM = 3$, what is the length of $BC$?


 


                    \begin{figure}[H]
                        \centering
                        \includegraphics[width=0.5\textwidth]{Problem/AHSME_I_1975_20_Problem_Figure_1.png}
                    \end{figure}
                    


 $\textbf{(A)}\ 2\sqrt{26} \qquad  \textbf{(B)}\ 2\sqrt{31} \qquad  \textbf{(C)}\ 9 \qquad  \textbf{(D)}\ 4+2\sqrt{13} \qquad\\ \textbf{(E)}\ \text{not enough information given to solve the problem}$


 



 

	\item (\textit{AHSME 1975}) Suppose $f(x)$ is defined for all real numbers $x$; $f(x)>0$ for all $x$, and $f(a)f(b)=f(a+b)$ for all $a$ and $b$. 
Which of the following statements is true? 


 $\text{I. } f(0)=1$


 $\text{II. } f(-a)=1/f(a) \text{ for all } a$


 $\text{III. } f(a) = \sqrt[3]{f(3a)} \text{ for all } a$


 $\text{IV. } f(b)>f(a) \text{ if } b>a$


 

$\textbf{(A)}\ \text{III and IV only} \qquad  \textbf{(B)}\ \text{I, III, and IV only} \qquad  \textbf{(C)}\ \text{I, II, and IV only} \qquad  \textbf{(D)}\ \text{I, II, and III only} \qquad  \textbf{(E)}\ \text{All are true.}$


 

	\item (\textit{AHSME 1975}) If $p$ and $q$ are primes and $x^2 - px + q = 0$ has distinct positive integral roots, 
then which of the following statements are true? 


 $I.\ \text{The difference of the roots is odd.} \\ II.\ \text{At least one root is prime.} \\ III.\ p^2-q\ \text{is prime}. \\ IV.\ p+q\ \text{is prime}$


 

$\textbf{(A)}\ \text{I only} \qquad  \textbf{(B)}\ \text{II only} \qquad  \textbf{(C)}\ \text{II and III only} \qquad  \textbf{(D)}\ \text{I, II and IV only} \qquad  \textbf{(E)}\ \text{All are true.}$


 

	\item (\textit{AHSME 1975}) In the adjoining figure $AB$ and $BC$ are adjacent sides of square $ABCD$; $M$ is the midpoint of $AB$; 
$N$ is the midpoint of $BC$; and $AN$ and $CM$ intersect at $O$. The ratio of the area of $AOCD$ to the area of $ABCD$ is


 
                    \begin{figure}[H]
                        \centering
                        \includegraphics[width=0.5\textwidth]{Problem/AHSME_I_1975_23_Problem_Figure_1.png}
                    \end{figure}
                    


 $\textbf{(A)}\ \frac{5}{6}\qquad \textbf{(B)}\ \frac{3}{4}\qquad \textbf{(C)}\ \frac{2}{3}\qquad \textbf{(D)}\ \frac{\sqrt{3}}{2}\qquad \textbf{(E)}\ \frac{(\sqrt{3}-1)}{2}$


 

	\item (\textit{AHSME 1975}) In $\triangle ABC$, $\measuredangle C=\theta$ and $\measuredangle B=2\theta$, where $0^{\circ} <\theta < 60^{\circ}$. 
The circle with center $A$ and radius $AB$ intersects $AC$ at $D$ and intersects $BC$, 
extended if necessary, at $B$ and at $E$ ($E$ may coincide with $B$). Then $EC=AD$


 $\textbf{(A)}\ \text{for no values of}\ \theta \qquad \textbf{(B)}\ \text{only if}\ \theta=45^{\circ} \qquad \textbf{(C)}\ \text{only if}\ 0^{\circ}<\theta\le 45^{\circ}\\ \qquad\ \textbf{(D)}\ \text{only if}\ 45^{\circ}\le\theta < 60^{\circ}\qquad \textbf{(E)}\ \text{for all}\ \theta\ \text{such that}\ 0^{\circ}<\theta < 60^{\circ}$


 

	\item (\textit{AHSME 1975}) A woman, her brother, her son and her daughter are chess players (all relations by birth). The worst player's twin (who is one of the four players) and the best player are of opposite sex. The worst player and the best player are the same age. Who is the worst player?


 $\textbf{(A)}\ \text{the woman} \qquad \textbf{(B)}\ \text{her son} \qquad \textbf{(C)}\ \text{her brother} \qquad \textbf{(D)}\ \text{her daughter}\\ \qquad \textbf{(E)}\ \text{No solution is consistent with the given information}$


 

	\item (\textit{AHSME 1975}) In acute $\triangle ABC$ the bisector of $\measuredangle A$ meets side $BC$ at $D$. The circle with center $B$ and radius $BD$ 
intersects side $AB$ at $M$; and the circle with center $C$ and radius $CD$ intersects side $AC$ at $N$. Then it is always true that


 $\textbf{(A)}\ \measuredangle CND+\measuredangle BMD-\measuredangle DAC =120^{\circ} \qquad \textbf{(B)}\ AMDN\ \text{is a trapezoid}\qquad \textbf{(C)}\ BC\ \text{is parallel to}\ MN\\ \qquad \textbf{(D)}\ AM-AN=\frac{3(DB-DC)}{2}\qquad \textbf{(E)}\ AB-AC=\frac{3(DB-DC)}{2}$


 

	\item (\textit{AHSME 1975}) If $p, q$ and $r$ are distinct roots of $x^3-x^2+x-2=0$, then $p^3+q^3+r^3$ equals


 $\textbf{(A)}\ -1 \qquad \textbf{(B)}\ 1 \qquad \textbf{(C)}\ 3 \qquad \textbf{(D)}\ 5 \qquad \textbf{(E)}\ \text{none of these}$


 

	\item (\textit{AHSME 1975}) In $\triangle ABC$ shown in the adjoining figure, $M$ is the midpoint of side $BC, AB=12$ and $AC=16$. Points $E$ and $F$ are taken on $AC$ 
and $AB$, respectively, and lines $EF$ and $AM$ intersect at $G$. If $AE=2AF$ then $\frac{EG}{GF}$ equals


 
                    \begin{figure}[H]
                        \centering
                        \includegraphics[width=0.5\textwidth]{Problem/AHSME_I_1975_28_Problem_Figure_1.png}
                    \end{figure}
                    


 $\textbf{(A)}\ \frac{3}{2} \qquad \textbf{(B)}\ \frac{4}{3} \qquad \textbf{(C)}\ \frac{5}{4} \qquad \textbf{(D)}\ \frac{6}{5}\\ \qquad \textbf{(E)}\ \text{not enough information to solve the problem}$


 

	\item (\textit{AHSME 1975}) What is the smallest integer larger than $(\sqrt{3}+\sqrt{2})^6$?


 $\textbf{(A)}\ 972 \qquad \textbf{(B)}\ 971 \qquad \textbf{(C)}\ 970 \qquad \textbf{(D)}\ 969 \qquad \textbf{(E)}\ 968$


 

	\item (\textit{AHSME 1975}) Let $x=\cos 36^{\circ} - \cos 72^{\circ}$. Then $x$ equals


 $\textbf{(A)}\ \frac{1}{3}\qquad \textbf{(B)}\ \frac{1}{2} \qquad \textbf{(C)}\ 3-\sqrt{6} \qquad \textbf{(D)}\ 2\sqrt{3}-3\qquad \textbf{(E)}\ \text{none of these}$


 

\end{enumerate}